\scanpage{16}
\mspage{14}
3.1\deleted{5}4. Soit \(\pi'\) un sous-groupe d'indice fini de \(\pi\). Il correspond donc à une extension séparable finie \(K'\) de \(K\), d'où un normalisé \(V'\) de \(V\) dans \(K'\), qui est un anneau de val. discrète hensélien. Si \(k'\) est son corps résiduel, \(k'/k\) est une extension finie, donc \(\pi'_0\unsure{\mathrel{-}|\longrightarrow}\pi'\) \deleted{\unsure{est|a}} \inserted{induit} un \deleted{\unsure{isom|isomes}} de \(\pi'_0\)\inserted{\(=\mathrm{Gal}(\bar k'/k')\)} sur un sous groupe ouvert de \(\pi_0=\mathrm{Gal}(\bar k/k)\). \deleted{D'autre part} \inserted{Il s'ensuit en particulier que \(\pi'_0\mathrel{-}\pi'\) est surjectif, i.e. que} \deleted{on voit que} \inserted{l'on a}
\[
 I'=I\cap\pi' \tag*{(3.1\deleted{5}\unsure{4|5}.1)}
\]

Notons que d'après le théor. général appliqué à \(V'\), on aura pour \unsure{tt|tout} \inserted{nb premier} \(\ell\ne p\)
\[
 [I'/(I',I')]^{\deleted{\unsure{(\ell)|\ell}}}(\ell)\simeq\mathbb Z_\ell(1) \tag*{(3.\unsure{14|15}.2)}
\]
(isom canonique). \deleted{D'autre part}, l'hom \(I'_t\mathrel{-}I_t\) s'identifie, sur \(I_t\), à la multiplication par \(e\) \inserted{comme dessus}, \deleted{qui} lorsque \(k'/k\) est séparable, \deleted{p.ex. que} \inserted{i.e.} l'extension résiduelle de \(\widetilde V'/\widetilde V\) \deleted{séparable} \inserted{triviale}, \deleted{\unsure{donc|dans|Alors}} p. ex. si \(k\) est parfait, alors on peut interpréter \(e\) en termes purement ``\unsure{gruppentheoretische|gruppentheoretischer}'' comme
\[
 e=[I:I']=\mathrm{card}\ I/I'.
\]
Sinon, il faut diviser par le facteur inséparable \(p^s\) de \(k'/k\).

\deleted{3.16. Il \unsure{reste à noter autrement|est intéressant}, lorsque \unsure{le corps|le fait} \(\unsure{I'|\pi'}\) invariant dans \(\pi\) \unsure{est invariant|est quotient} de \(I\), \(\pi'/\pi\) opère sur}

3.15. Lorsque \(\pi'\) est invariant dans \(\pi\) \inserted{(i.e. \(K'/K\) galoisienne)}, il en est de \(\overline{\mathrm m}\) de \deleted{\(\pi'_0\) dans \(\pi_0\), puisque (on ne sait pas \(k'/k\) est une ext. \unsure{quasi-|pure-}galoisienne), donc \unsure{l'hom|l'}} \(I'=I\cap\pi'\) \inserted{dans \(\pi\)} (puisque \(I\) est invariant), \deleted{de \(\unsure{\pi/\pi'|\pi'/\pi}\)} \deleted{\unsure{invariant|invariance}} \(\pi\unsure{?|\ne}\) opère sur \(I'/(I',I')^\wedge\) et par suite sur

% archive sheet number: 15
